Computational mouse-brain studies combine evidence with distinct meanings. Reproducible execution or accurate prediction may support a bounded result while leaving mechanism, causality, or digital-twin scope unauthorized. MouseClaimBench is an integrity-aware knowledge-based system whose non-compensatory contracts bind domain evidence, structural admissibility, provenance integrity, and complete deficits. For policy traceability, profile v2 provides a machine-readable acquisition ledger linking 10 claims, 22 predicate contracts, 60 claim--evidence relations, and 124 observation slots to sources and relation-specific rationales. The profile remains author-defined rather than independently content-validated. Python and \ac{shacl} agree on all 5,497 structural contract cases, while an independently implemented \ac{asp} path agrees on a deterministically selected subset of 262 cases. A further 55,031 property checks show no violation. Across 221 profiles and 14 fixed decisions, one removal expands one authorization and 52 additions produce 71 contractions. Counterfactual tests confirm 9,839 sufficient repairs and 51,480 individual-necessity checks. The complete integrity gate blocks all 2,550 non-empty compositions of eight declared attacks, while 20 internally coherent forgery controls intentionally escape because no external trust anchor is available. Two pre-access frozen \ac{dandi} protocols preserve one availability refusal and one bounded predictive authorization across 32 mice. A dependence-aware \ac{microns} application authorizes only a local observational structure--function association. These results establish traceable policy execution and quantified policy dependence, not biological truth, human utility, expert consensus, or a complete digital mouse brain.
